% TOP_PROBLEM: {"id": 30003733, "title": "Low-Parameter Hyperelastic Models of Brain Tissue", "category_id": 16, "category": {"id": 16, "name": "physics", "display_name": "Mathematical Physics", "description": "Problems at the intersection of mathematics and physics.", "slug": "physics", "order_index": 16, "created_at": "2026-07-31T15:26:25.671Z"}, "status": "partially_solved", "statusQualification": "This is an experimentally constrained constitutive-model question. The report\\textquotesingle s observed shear/compression/tension trends and the small-to-moderate in-vivo strain regime are part of the target, not a request for a universal energy independent of tissue or data. ID 30003734 is a duplicate extraction."}
% FIRST_POSTED: 2026-10-11
% ENTRY_KIND: statement_only
% SUBMISSION: {"kind": "statement_only", "category": "top_problems", "problem_id": "30003733", "proposal_type": "established", "source": "The Mathematics of Mechanobiology and Cell Signaling (2018), Oberwolfach Reports, DOI 10.4171/owr/2018/8, pp. 433--506. Primary report accessed 2026-10-11; relevant PDF pages 12. https://doi.org/10.4171/owr/2018/8; original OWR report (catalogue-listed context; not independently consulted in this preparation). https://doi.org/10.4171/OWR/2018/8; Mihai et al., hyperelastic brain models (catalogue-listed context; not independently consulted in this preparation). https://doi.org/10.1016/j.jmps.2017.05.015; Filla et al., heterogeneous brain constitutive model (catalogue-listed context; not independently consulted in this preparation). https://doi.org/10.1016/j.jmbbm.2023.106271", "status_evidence": "Canonical source status partially_solved retained. Supplied literature triage dated 2026-08-21: Validated constitutive models now reproduce important brain-tissue responses under multiaxial loading, but no universally accepted first-principles model covers the tissue's regional, viscoelastic, and multiscale behavior. Primary-report access checked 2026-10-11; this historical report and inherited triage do not certify current resolution of every component.", "status_checked_on": "2026-10-11", "duplicate_check": "Compared canonical ID, title and normalized target with the complete local UnsolvedMath registry; normalized exact-text matches: none; title matches: none. Related or overlapping canonical records: [30003734]. Primary source context and the existing OWR statement audit used for scope and duplicate review; an exhaustive mathematical-equivalence/literature search was not performed.", "attribution": "Original proposer not named in the supplied canonical record; attribution remains to the identified Oberwolfach contribution/report, not its editors as invented proposers. Explanatory statement preparation and source/semantic audit assisted by OpenAI GPT-6 Codex; no human mathematical authorship or independent proof verification is claimed."}
% SOURCE_STATEMENT: "Determine a low-parameter family of isotropic hyperelastic strain-energy functions for brain tissue that captures its observed response under combined shear and axial compression or tension, remains suitable for small-to-moderate in-vivo strains, and can be calibrated and validated on multiaxial data."
% Preparation and final self-audit: CONTRIBUTING.md prompts C and D; source/audit artifacts retained in the local physics_statements_2026_10_11 batch.
% Imported from: physics_statements.tex; UnsolvedMath ID 30003733
% !TeX program = lualatex
\documentclass[11pt]{article}
\usepackage[margin=25mm]{geometry}
\usepackage{fontspec}
\setmainfont{Latin Modern Roman}
\usepackage{amsmath,amssymb,mathtools}
\usepackage{unicode-math}
\setmathfont{Latin Modern Math}
\usepackage{xurl,hyperref}
\hypersetup{colorlinks=true,urlcolor=blue}
\setcounter{secnumdepth}{0}
\setlength{\parindent}{0pt}
\setlength{\parskip}{5pt}
\setlength{\emergencystretch}{3em}
\newcommand{\sref}[2]{\hyperlink{source:#1:#2}{[S#2]}}
\newcommand{\cataloguescope}[1]{\paragraph{Recorded target.}#1}
\newcommand{\R}{\mathbb R}
\newcommand{\C}{\mathbb C}
\newcommand{\N}{\mathbb N}
\newcommand{\Z}{\mathbb Z}
\newcommand{\Q}{\mathbb Q}
\newcommand{\F}{\mathbb F}
\newcommand{\eps}{\varepsilon}
\begin{document}
% UnsolvedMath ID 30003733; source catalogue code OWR-15993-001
\setcounter{section}{30003732}
\section{Low-Parameter Hyperelastic Models of Brain Tissue}\label{30003733}
{\small UnsolvedMath ID 30003733 · Mathematical Physics}
\subsection{Definitions and mathematical statement}

\paragraph{Shared notation.}$\mathbb N=\{1,2,\ldots\}$; the standard number fields have their usual meanings. All additional conventions are stated locally.


\paragraph{Definitions and conventions.}
A hyperelastic material has strain energy $W(F)$ depending on deformation gradient $F$; stress follows by differentiating this energy. Objectivity and isotropy impose invariance under appropriate rotations of $F$. Shear and axial compression or tension probe different combinations of its invariants. A low-parameter family uses few calibratable constitutive constants; validation requires comparison with independent multiaxial data, not merely interpolation of one loading curve.

\paragraph{Target retained from the canonical source catalogue.}
\paragraph{Catalogue wording retained for traceability.}
Determine a low-parameter family of isotropic hyperelastic strain-energy functions for brain tissue that captures its observed response under combined shear and axial compression or tension, remains suitable for small-to-moderate in-vivo strains, and can be calibrated and validated on multiaxial data.

\paragraph{Source scope and qualification.}
This is an experimentally constrained constitutive-model question. The report\textquotesingle s observed shear/compression/tension trends and the small-to-moderate in-vivo strain regime are part of the target, not a request for a universal energy independent of tissue or data. ID 30003734 is a duplicate extraction.
\subsection{Short English statement}
Find a small, testable hyperelastic model that captures the unusual combined-load response of brain tissue.
\subsection{Sources}
\begin{itemize}
\item[\hypertarget{source:30003733:1}{S1}] ULAM.AI and the UnsolvedMath Contributors, UnsolvedMath, record 30003733 (OWR-15993-001). Catalogue curation under CC BY 4.0; definitions and English exposition added locally.
\url{https://huggingface.co/datasets/ulamai/UnsolvedMath}.
\item[\hypertarget{source:30003733:2}{S2}] The Mathematics of Mechanobiology and Cell Signaling (2018), Oberwolfach Reports, DOI 10.4171/owr/2018/8, pp. 433--506. Primary report accessed 2026-10-11; relevant PDF pages 12.
\url{https://doi.org/10.4171/owr/2018/8}.
\item[\hypertarget{source:30003733:3}{S3}] original OWR report (catalogue-listed context; not independently consulted in this preparation).
\url{https://doi.org/10.4171/OWR/2018/8}.
\item[\hypertarget{source:30003733:4}{S4}] Mihai et al., hyperelastic brain models (catalogue-listed context; not independently consulted in this preparation).
\url{https://doi.org/10.1016/j.jmps.2017.05.015}.
\item[\hypertarget{source:30003733:5}{S5}] Filla et al., heterogeneous brain constitutive model (catalogue-listed context; not independently consulted in this preparation).
\url{https://doi.org/10.1016/j.jmbbm.2023.106271}.
\end{itemize}
\label{end:30003733}
\end{document}
